% ==============================================================================
% BAB V: KESIMPULAN DAN SARAN
% ==============================================================================

\chapter{KESIMPULAN DAN SARAN}

\section{Kesimpulan}
Berdasarkan hasil perancangan, pengujian eksperimen empiris, dan analisis komparatif yang telah dilakukan terhadap NGINX dan HAProxy pada lingkungan kontainer Docker dengan kapasitas \textit{backend} heterogen, dapat ditarik kesimpulan sebagai berikut:
\begin{enumerate}
    \item Algoritma penjadwalan berbobot (\textit{Weighted Round Robin} dan \textit{Weighted Least Connections}) secara signifikan mengungguli algoritma tanpa bobot (\textit{Round Robin} dan \textit{Least Connections}) dalam kluster kontainer Docker heterogen. Algoritma WLC menghasilkan kinerja paling optimal dengan mencatatkan peningkatan \textit{throughput} hingga 60\% serta meredam latensi rata-rata hingga 41\% dibandingkan algoritma Round Robin konvensional.
    \item Pada karakteristik beban kerja intensif komputasi (\textit{CPU-bound}), ketidaksesuaian algoritma tanpa bobot memicu fenomena \textit{CPU throttling} parah pada kontainer berkapasitas rendah, yang berakibat pada lonjakan ekstrem \textit{tail latency} (p95 dan p99) hingga di atas 4.000 ms. Penerapan algoritma berbobot berhasil mengeliminasi degradasi latensi ekor tersebut hingga berada pada rentang stabil 1.490--1.610 ms.
    \item Dalam perbandingan arsitektural \textit{head-to-head}, HAProxy menunjukkan performa data-plane yang sedikit lebih unggul dan lebih konsisten dibandingkan NGINX, terutama dalam menstabilkan latensi ekor p99 pada beban kerja intensif I/O dengan konkurensi tinggi berkat efisiensi struktur data antrean \texttt{ebtree} dan arsitektur \textit{event-driven single-process}.
    \item Audit log distribusi trafik membuktikan bahwa algoritma WRR dan WLC pada kedua mesin penyeimbang beban mampu merealisasikan pembagian beban aktual yang sangat presisi terhadap rasio kapasitas kuota cgroups 1:2:4 dengan batas simpangan galat di bawah 1,5\%.
\end{enumerate}

\section{Saran}
Beberapa saran yang dapat diajukan untuk pengembangan penelitian selanjutnya antara lain:
\begin{enumerate}
    \item Mengembangkan mekanisme penyesuaian bobot secara dinamis (\textit{dynamic auto-weighting}) berbasis telemetri utilisasi metrik cgroups secara \textit{real-time}, sehingga bobot kontainer dapat menyesuaikan diri secara otomatis saat terjadi fluktuasi beban eksternal.
    \item Memperluas skenario pengujian pada lingkungan orkestrasi kontainer tingkat lanjut seperti Kubernetes dengan mengevaluasi perbandingan kinerja antara Ingress Controller berbasis NGINX, HAProxy, Envoy, dan Traefik.
    \item Mengintegrasikan pengujian protokol HTTP/3 (QUIC) untuk mengevaluasi dampak protokol transport berbasis UDP terhadap kinerja penyeimbang beban pada kluster heterogen.
\end{enumerate}
